% ============================================================
% ch11.tex —— 第 11 章 极大表示（Maximal representations）
% 底本：Springer LNM 2386，印刷 pp.111–128（PDF p151–p168）
% 页锚点：正文中形如「原文 p.xx」的注释标记页界，xx 为印刷页码
% ============================================================

\chapter{极大表示}

% [原文 p.111]

\ChapterAbstract{$G$-特征簇的连通分支由取值于 $\pi_1(G)$ 的初等拓扑不变量加以区分。对埃尔米特李群而言，这个不变量就是 Toledo 不变量，它的绝对值被一个 Milnor--Wood 型不等式所界。在管型的埃尔米特群的情形，Toledo 不变量取到最大值的那些连通分支称为极大分支。它们构成另一类高阶泰希米勒空间。本章我们讨论使人们能够显式研究它们的坐标系。}

\section{极大表示}

管型埃尔米特李群构成对高阶泰希米勒理论极为重要的一类半单李群。与分裂实李群类似，管型埃尔米特李群具有某种正性概念，我们将在本节以及第 12.4 节中讨论它。此外，粗略地说，管型埃尔米特李群可以看作“大” $\SL_2(\R)$——更准确地说“大” $\Sp_2(\R)$——李群。这意味着，可以把它们当作带有某种正性结构的、非交换代数上的辛李群来处理。关于管型埃尔米特李群的这一视角的更多细节，我们请读者参阅 [13]。

将曲面群的极大表示映入管型埃尔米特李群，所得的空间提供了高阶泰希米勒空间的又一例子。在本节中，我们定义埃尔米特李群与极大表示，并讨论它们的若干性质。

\subsection{埃尔米特对称空间与埃尔米特李群}

设 $M$ 是连通复流形，$g$ 是 $M$ 上的埃尔米特度量。

% [原文 p.112]

\begin{definition}[11.1]
$(M,g)$ 称为埃尔米特对称空间（Hermitian symmetric space），如果每个点 $p\in M$ 都是某个全纯对合等距 $s_p\colon M\to M$ 的孤立不动点。
\end{definition}

\begin{theorem}[11.2]（Harish--Chandra 嵌入定理 {[195, 定理 7.1]}）
每个非紧埃尔米特对称空间都双全纯于复向量空间中的一个开有界域。
\end{theorem}

\begin{definition}[11.3]
埃尔米特对称空间称为管型的（of tube type），如果它双全纯于形如 $V+\ii\Omega$ 的区域，其中 $V$ 是实向量空间、$\Omega\subset V$ 是开的真凸锥。
\end{definition}

% [Tutor Note: 用 Def. 11.3 核对本节的主要例子：Sp_{2n}(R)/K ≅ S = Sym_n(R) + i Sym_n^+(R)，这里 V = Sym_n(R) 为实向量空间，Ω = Sym_n^+(R)（正定对称矩阵全体）是开的真凸锥，故 Siegel 上半空间恰为管型域，Sp_{2n}(R) 确实属于 Fact 11.5 的管型名单。两行名单的差别只在第 2、3 个群：SU(n,m) 仅当 m = n 时管型（故第二行写 SU(n,n)），SO^*(2n) 仅当 n 为偶数时管型（故第二行写 SO^*(4n)）。]

设 $G$ 是连通非紧半单李群，其中心有限（或更一般地，紧），$K$ 是 $G$ 的极大紧子群。

\begin{definition}[11.4]
群 $G$ 称为埃尔米特李群（Hermitian Lie group），如果对称空间 $G/K$ 是埃尔米特的。若对称空间 $G/K$ 是管型的，则称 $G$ 为管型埃尔米特李群。
\end{definition}

\begin{fact}[11.5]
埃尔米特李群已被分类。在覆叠的意义下，它们是：
\[
\Sp_{2n}(\R),\quad \SU(n,m),\quad \SO^{*}(2n),\quad \SO_0(2,n+2),\quad E_{7(-25)},\quad E_{6(-14)}
\]
其中 $n,m\in\N$。前四个群称为经典的，后两个群称为例外的。而群
\[
\Sp_{2n}(\R),\quad \SU(n,n),\quad \SO^{*}(4n),\quad \SO_0(2,n+2),\quad E_{7(-25)}
\]
是管型的。
\end{fact}

本节我们的主要例子将是实辛群 $\Sp_{2n}(\R)$。它的极大紧子群
\[
K=\left\{\begin{pmatrix} u_1 & u_2\\ -u_2 & u_1\end{pmatrix}\ \middle|\ u_1+\ii u_2\in\U(n)\right\}
\]
同构于 $\U(n)$。对称空间 $\Sp_{2n}(\R)/K$ 可以等同于 Siegel 上半空间
\[
\mathcal{S}:=\left\{a+\ii b\in\Sym_n(\C)\ \middle|\ a\in\Sym_n(\R),\ b\in\Sym_n^{+}(\R)\right\}.
\]
群 $G$ 通过广义 Möbius 变换作用在 $\mathcal{S}$ 上：
\[
\begin{pmatrix} a & b\\ c & d\end{pmatrix}.z=(az+b)(cz+d)^{-1},\qquad
\text{其中 } z\in\mathcal{S},\ \begin{pmatrix} a & b\\ c & d\end{pmatrix}\in\Sp_{2n}(\R).
\]

\subsection{拉格朗日格拉斯曼流形}

% [原文 p.113]

考虑辛向量空间 $(\R^{2n},\omega)$，其中 $\omega$ 是 $\R^{2n}$ 上的标准辛形式，即
\begin{equation}
\omega(x,y)=x^{t}\begin{pmatrix} 0 & \mathrm{Id}\\ -\mathrm{Id} & 0\end{pmatrix}y, \tag{11.1}
\end{equation}
其中 $x$ 与 $y$ 属于 $\R^{2n}$。

若 $\R^{2n}$ 的一组基使得 $\omega$ 在这组基下的表达式具有 (11.1) 的形式，则称它为辛基（symplectic basis）（因此标准基是辛基）。我们通常把辛基写成一对 $(\mathbf{e},\mathbf{f})$，其中 $\mathbf{e}=(e_1,\dots,e_n)$、$\mathbf{f}=(f_1,\dots,f_n)$；于是对一切 $i,j$ 都有 $\omega(e_i,f_j)=\delta_{ij}$。更一般地，给定两族 $\mathbf{v}=(v_1,\dots,v_n)$ 与 $\mathbf{w}=(w_1,\dots,w_m)$，我们用 $\omega(\mathbf{v},\mathbf{w})$ 表示系数为 $\omega(v_i,w_j)$ 的 $m\times n$ 矩阵；此时有 $\omega(\mathbf{w},\mathbf{v})=-\omega(\mathbf{v},\mathbf{w})^{t}$。

实辛群可以看作形式 $\omega$ 的等距群：
\[
\Sp_{2n}(\R)=\left\{g\in\GL_{2n}(\R)\mid g^{t}\omega g=\omega\right\}.
\]
它的伴随形式 $\mathrm{PSp}_{2n}(\R):=\Sp_{2n}(\R)/\{\pm\mathrm{Id}\}$ 称为射影实辛群。

\begin{definition}[11.6]
$\R^{2n}$ 的子空间 $L$ 称为拉格朗日的（Lagrangian），如果 $\dim(L)=n$ 且对一切 $u,v\in L$ 有 $\omega(u,v)=0$。$(\R^{2n},\omega)$ 的全部拉格朗日子空间构成的空间称为拉格朗日格拉斯曼流形，记作 $\Lag_n$。两个拉格朗日子空间 $L_1,L_2$ 称为横截的（transverse），如果它们的交为零。
\end{definition}

\begin{remark}[11.7]
\begin{enumerate}[label=\arabic*.,leftmargin=*]
\item 回忆：对任意 $n$ 维实向量空间 $L$，群 $\GL_n(\R)$ 从右作用在 $L$ 的基构成的空间上：若 $\mathbf{v}=(v_1,\dots,v_n)$ 是 $L$ 的一组基，而 $g=(g_{i,j})_{i,j\in\{1,\dots,n\}}\in\GL_n(\R)$，则 $\mathbf{w}:=\mathbf{v}g$ 是由
\[
w_j=\sum_{i=1}^{n}v_i g_{i,j},\qquad \text{对一切 } j\in\{1,\dots,n\}
\]
定义的族 $(w_1,\dots,w_n)$。这正是上面提到的主丛结构背后的单传递作用。
\item 在维数为 $2n$ 的辛向量空间的情形，我们由此得到 $\Sp_{2n}(\R)$ 在辛基构成的空间上的单传递作用。
\item 当 $\mathbf{v}$ 是向量空间 $V$ 中 $n$ 个元素构成的族（不必线性无关，也不必生成 $V$），而 $g$ 是 $m\times n$ 矩阵时，我们也将使用记号 $\mathbf{v}\cdot g$（有时写作 $\mathbf{v}g$）。
\item 态射 $\psi\colon L\to L'$ 与它关于 $L$ 的基 $\mathbf{v}$、$L'$ 的基 $\mathbf{v}'$ 的矩阵 $g$ 由众所周知的公式相联系：$\psi(\mathbf{v})=\mathbf{v}'g$。
\item 对两个族 $\mathbf{v}=(v_1,\dots,v_n)$ 与 $\mathbf{w}=(w_1,\dots,w_n)$，族 $(v_1+w_1,\dots,v_n+w_n)$ 记作 $\mathbf{v}+\mathbf{w}$。
\end{enumerate}
\end{remark}

群 $\Sp_{2n}(\R)$ 在 $\Lag_n$ 上有自然的左作用，由
\[
g\cdot L:=\{g(x)\}_{x\in L}
\]
给出。这个作用是传递的，因此 $\Lag_n$ 是辛群下的齐性空间。

设 $(e_0,f_0)$ 是 $\R^{2n}$ 的标准辛基。拉格朗日子空间 $L_0=\operatorname{Span}(e_0)$ 与 $L_0^{\mathrm{opp}}=\operatorname{Span}(f_0)$ 称为标准的（standard）。注意它们是横截的。考虑稳定化子
\[
P=\Stab_{\Sp_{2n}(\R)}(L_0)
=\left\{\begin{pmatrix} A & AB\\ 0 & (A^{t})^{-1}\end{pmatrix}\ \middle|\ A\in\GL_n(\R),\ B\in\Sym_n(\R)\right\},
\]
% [Tutor Note: 这里右上块印作 AB（不是 B），并非笔误：设 g = (A X; 0 D)，由 g^t ω g = ω（ω = (0 Id; -Id 0)）的右上块条件得 A^T D = Id，即 D = (A^t)^{−1}；右下块条件给出 X^T D = D^T X，等价于 A^{−1}X ∈ Sym_n(R)。取 X = AB、B ∈ Sym_n(R) 恰好满足 A^{−1}X = B。同理下一页 P^opp 的左下块印作 BA（这来自 A^T C = C^T A，取 C = BA、B ∈ Sym_n(R) 即可）。若误以为应写成 B，反而不满足辛条件，可据此自查。]

% [原文 p.114]

\[
U=\Stab_{\Sp_{2n}(\R)}(\mathbf{e}_0)
=\left\{\begin{pmatrix} I_n & B\\ 0 & I_n\end{pmatrix}\ \middle|\ B\in\Sym_n(\R)\right\},
\]
\[
P^{\mathrm{opp}}=\Stab_{\Sp_{2n}(\R)}(L_0^{\mathrm{opp}})
=\left\{\begin{pmatrix} A & 0\\ BA & (A^{t})^{-1}\end{pmatrix}\ \middle|\ A\in\GL_n(\R),\ B\in\Sym_n(\R)\right\},
\]
\[
U^{\mathrm{opp}}=\Stab_{\Sp_{2n}(\R)}(\mathbf{f}_0)
=\left\{\begin{pmatrix} I_n & 0\\ B & I_n\end{pmatrix}\ \middle|\ B\in\Sym_n(\R)\right\},
\]
其中 $I_n$ 是 $n\times n$ 单位矩阵。

群 $P$ 与 $P^{\mathrm{opp}}$ 是 $\Sp_{2n}(\R)$ 的互为对顶的极大抛物子群，而 $U\subset P$ 与 $U^{\mathrm{opp}}\subset P^{\mathrm{opp}}$ 是它们的幂幺根。此外，我们有幂幺根的李代数：
\[
\mathfrak{u}=\operatorname{Lie}(U)
=\left\{\begin{pmatrix} 0 & B\\ 0 & 0\end{pmatrix}\ \middle|\ B\in\Sym_n(\R)\right\},
\]
\[
\mathfrak{u}^{\mathrm{opp}}=\operatorname{Lie}(U^{\mathrm{opp}})
=\left\{\begin{pmatrix} 0 & 0\\ B & 0\end{pmatrix}\ \middle|\ B\in\Sym_n(\R)\right\},
\]
并且 $P$ 与 $P^{\mathrm{opp}}$ 的交称为 $\Sp_{2n}(\R)$ 的 Levi 子群：
\[
L:=P\cap P^{\mathrm{opp}}
=\left\{\begin{pmatrix} A & 0\\ 0 & (A^{t})^{-1}\end{pmatrix}\ \middle|\ A\in\GL_n(\R)\right\}.
\]
Levi 子群通过伴随作用（在我们的情形也就是共轭）作用在 $\mathfrak{u}$ 上，并保持下述尖锐开凸锥
\[
\mathfrak{u}^{+}=\left\{\begin{pmatrix} 0 & B\\ 0 & 0\end{pmatrix}\ \middle|\ B\in\Sym_n^{+}(\R)\right\}.
\]

由于 $\Sp_{2n}(\R)$ 在 $\Lag_n$ 上的作用是传递的，拉格朗日格拉斯曼流形是光滑流形，微分同胚于齐性空间 $\Sp_{2n}(\R)/P$。$\Sp_{2n}(\R)$ 在 $\Lag_n$ 上的作用不是有效的，其核为 $\{\pm\mathrm{Id}\}$。$\Lag_n$ 的对称群是射影辛群 $\mathrm{PSp}_{2n}(\R)$。

\subsection{拉格朗日的构形}

对每个整数 $d\ge 2$，用 $\Conf^{d}(\Lag_n)$ 表示 $d$ 元拉格朗日组的模空间，即 $(\Lag_n)^{d}$ 模去 $\Sp_{2n}(\R)$ 的对角作用所得的商。对称群 $\mathfrak{S}_d$ 在 $(\Lag_n)^{d}$ 上的自然作用下降为 $\Conf^{d}(\Lag_n)$ 上的作用。

我们将特别关注 $\Conf^{2}(\Lag_n)$、$\Conf^{3}(\Lag_n)$ 及它们的某些子空间。我们用 $\Conf^{d*}(\Lag_n)$ 表示 $d$ 个两两横截的拉格朗日组的构形空间。子空间 $\Conf^{d*}(\Lag_n)\subset\Conf^{d}(\Lag_n)$ 在一切置换下不变。一个容易的线性代数练习表明，实辛群在 $\Conf^{2*}(\Lag_n)$ 上传递地作用，且 $\Conf^{2*}(\Lag_n)$ 的任一元素的稳定化子在 $\Sp_{2n}(\R)$ 中共轭于 Levi 子群。

% [原文 p.115]

\subsection{Maslov 指标}

本节我们回顾 Maslov 指标的诸性质。

设 $L_1,M,L_2$ 是三个两两横截的拉格朗日子空间。存在唯一的线性映射 $M_{L_1\to L_2}\colon L_1\to L_2$，使得
\[
M=\left\{v\in\R^{2n}\ \middle|\ \text{存在 } e\in L_1,\ \text{使得 } v=e+M_{L_1\to L_2}(e)\right\}.
\]
在不引起混淆时，就把它简记作 $M$。利用辛形式 $\omega$，可以如下在 $L_1$ 上定义一个双线性形式 $\beta$：对 $v_1,v_2\in L_1$，
\[
\beta(v_1,v_2):=\omega(v_1,M(v_2)).
\]

\begin{definition}[11.8]
双线性形式 $\beta$ 称为 Maslov 形式（Maslov form），记作 $[L_1,M,L_2]$。
\end{definition}

下面的命题是一个容易的线性代数练习：

\begin{proposition}[11.9]
Maslov 形式 $[L_1,M,L_2]$ 是对称且非退化的。
\end{proposition}

\begin{remark}[11.10]
设 $\mathbf{e}$ 是 $L_1$ 的基，$\mathbf{f}$ 是 $L_2$ 的基，使得 $(\mathbf{e},\mathbf{f})$ 是辛基。那么 $M$ 在这些基下的矩阵 $[M]_{\mathbf{e},\mathbf{f}}$ 与矩阵 $[\beta]_{\mathbf{e}}$ 相等：$[M]_{\mathbf{e},\mathbf{f}}=[\beta]_{\mathbf{e}}$。
\end{remark}

我们用
\[
\sgn(\beta):=p-q
\]
表示 $\beta$ 的符号差，其中 $p$ 是 $L_1$ 上使 $\beta$ 正定的极大子空间的维数，$q$ 是 $L_1$ 上使 $\beta$ 负定的极大子空间的维数。它们满足 $p+q=n$，从而 $\sgn(\beta)=n \bmod 2$。

\begin{definition}[11.11]
拉格朗日三元组 $(L_1,M,L_2)$ 的 Maslov 指标是符号差 $\sgn([L_1,M,L_2])$，记作 $\mu_n(L_1,M,L_2)$。三元组 $(L_1,M,L_2)$ 称为极大的，如果它的 Maslov 指标等于 $n$。三元组 $(L_0,L^{+},L_0^{\mathrm{opp}})$（其中 $L^{+}$ 是由基 $\mathbf{e}_0+\mathbf{f}_0$ 生成的拉格朗日子空间）称为标准极大三元组。
\end{definition}

对 $n=1$，三个拉格朗日子空间 $(L_1,M,L_2)$ 是圆周 $\R\mathrm{P}^{1}$ 中三个互不相同的点。若这三点的循环次序为正，则 Maslov 指标为 $1$，否则为 $-1$。

% [原文 p.116]

\begin{remark}[11.12]
存在一个稍为一般的 Maslov 指标定义，它对任意拉格朗日三元组都适用，而不只对两两横截的三元组。它可以定义为二次形式
\[
\gamma\colon L_1\oplus M\oplus L_2\longrightarrow\R,\qquad
(v,w,x)\longmapsto\omega(v,w)+\omega(w,x)+\omega(x,v)
\]
的符号差 $\sgn(\gamma)$。当三元组两两横截时，这两个定义一致。
\end{remark}

% [Tutor Note: 核对 Maslov 指标的取值：Def. 11.11 用符号差定义，而 β 是 n 维空间 L_1 上的非退化双线性形式，故 p+q=n，于是 sgn(β) ≡ n (mod 2) 且 |sgn(β)| ≤ n，这正好与 Prop. 11.13 第一条的取值范围 {−n,−n+2,…,n} 相符（n 与 −n 同奇偶，故该集合就是所有与 n 同奇偶的整数）。n=1 时三点在 RP^1 上，μ=1 当且仅当循环有序，恰给出极值 ±1，可用来自查符号约定。]

\begin{proposition}[11.13]（Maslov 指标的性质）
Maslov 指标
\begin{itemize}
\item 的取值范围是 $\{-n,-n+2,\dots,n\}$；
\item 在 $\Sp_{2n}(\R)$ 于 $\Lag_n^{3}$ 上的作用下不变；
\item 是反对称的，因而也是循环不变的；
\item 满足上闭链关系，即对一切 $L_1,L_2,L_3,L_4\in\Lag_n$，成立
\[
\mu_n(L_1,L_2,L_3)-\mu_n(L_1,L_2,L_4)+\mu_n(L_1,L_3,L_4)-\mu_n(L_2,L_3,L_4)=0.
\]
\end{itemize}
群 $\Sp_{2n}(\R)$ 在具有相同 Maslov 指标的两两横截的拉格朗日三元组的集合上传递地作用。三元组 $(L_1,L_2,L_3)$ 是极大的，当且仅当存在元素 $g\in\Sp_{2n}(\R)$ 与元素 $u\in\mathfrak{u}^{+}$，使得 $g(L_1,L_2,L_3)=(L_0,\exp(u)L_0^{\mathrm{opp}},L_0^{\mathrm{opp}})$。
\end{proposition}

\begin{proposition}[11.14]（对拉格朗日极大三元组的作用）
群 $\Sp_{2n}(\R)$ 在拉格朗日极大三元组上传递地作用，且
\[
\mathcal{K}=\left\{\begin{pmatrix} u & 0\\ 0 & u\end{pmatrix}\ \middle|\ u\in\mathrm{O}(n)\right\}\cong\mathrm{O}(n)
\]
是标准极大三元组的稳定化子。
\end{proposition}

\begin{definition}[11.15]
拉格朗日四元组 $(L_1,L_2,L_3,L_4)$ 称为正的（positive），如果三元组 $(L_1,L_2,L_3)$ 与 $(L_1,L_3,L_4)$ 都是极大的。
\end{definition}

\begin{proposition}[11.16]（对正拉格朗日四元组的作用）
设 $(L_1,L_2,L_3,L_4)$ 是正四元组。则存在元素 $g\in\Sp_{2n}(\R)$ 与 $b\in\Sym_n^{+}(\R)$，使得 $g(L_1,L_2,L_3,L_4)=(L_0,L^{+},L_0^{\mathrm{opp}},L(b))$，其中 $L(b)$ 是由基 $\mathbf{f}_0-\mathbf{e}_0 b$ 生成的拉格朗日子空间。元素 $g$ 在左乘 $\mathcal{K}$ 中元素的意义下唯一。元素 $b$ 由四元组 $(L_1,L_2,L_3,L_4)$ 在 $\mathrm{O}(n)$ 中元素共轭的意义下唯一确定。
\end{proposition}

\begin{remark}[11.17]
在上一命题的记号下，$b$ 的特征值由四元组 $(L_1,L_2,L_3,L_4)$ 完全确定，并且全为正。$b$ 的谱有时称为正四元组 $(L_1,L_2,L_3,L_4)$ 的交比（cross ratio）。
\end{remark}

% [原文 p.117]

\subsection{映入管型埃尔米特李群的极大表示}

设 $S$ 是有限型的定向曲面。同态 $\rho\colon\pi_1(S)\to\Sp_{2n}(\R)$ 的一个重要不变量是 Toledo 不变量 $T_\rho$，它在 [70] 中借助带穿孔或边界的曲面的有界上同调引入，在 [69] 中则针对闭曲面引入。需要强调的是，在带穿孔或边界的曲面的情形，Toledo 数依赖于拓扑曲面 $S$，而不只依赖于它的基本群。它是一个实数，并满足 Milnor--Wood 型不等式：
\[
-n|\chi(S)|\le T_\rho\le n|\chi(S)|.
\]
在闭曲面的情形，这个数还是整数。此外，$T_\rho$ 在 $\Sp_{2n}(\R)$ 通过共轭作用于 $\Hom(\pi_1(S),\Sp_{2n}(\R))$ 之下不变，因此它对 $\mathcal{E}(S_{g,n},\Sp_{2n}(\R))$ 中的表示也有良好定义。

\begin{definition}[11.18]
表示 $\rho\in\mathcal{E}(S,\Sp_{2n}(\R))$ 称为极大的（maximal），如果 $T_\rho=n|\chi(S)|$。$\mathcal{E}(S,\Sp_{2n}(\R))$ 中所有极大表示构成的子空间记作 $\mathcal{M}(S,\Sp_{2n}(\R))$。
\end{definition}

极大表示具有重要的几何性质，我们在下一定理中陈述它们。这些性质已在 [70, 69, 14] 中证明。

\begin{theorem}[11.19]
设 $S$ 是有限型的定向曲面。
\begin{itemize}
\item 每个极大表示 $\rho$ 都是约化的，即 $\mathcal{M}(S,\Sp_{2n}(\R))\subset\mathcal{R}(S,\Sp_{2n}(\R))$。
\item 每个极大表示 $\rho$ 都是离散且忠实的，即 $\rho$ 的核为零，且 $\rho$ 的像在 $\Sp_{2n}(\R)$ 中离散。
\item 空间 $\mathcal{M}(S,\Sp_{2n}(\R))$ 是 Hausdorff 的。
\item 若 $S$ 是闭的，则 $\mathcal{M}(S,\Sp_{2n}(\R))$ 是 $\mathcal{R}(S,\Sp_{2n}(\R))$ 的开且闭的子集，即它是 $\mathcal{R}(S,\Sp_{2n}(\R))$ 的一些连通分支的并。
\item 同态 $\rho$ 是极大的，当且仅当它容许一个连续的极大边界映射，即存在映射 $\xi\colon\partial\pi_1(S)\to\Lag_n$，它是 $\rho$-等变的，并且是极大的，意思是：对每个正三元组 $(x_1,x_2,x_3)\in\partial\pi_1(S)^{3}$，像 $(\xi(x_1),\xi(x_2),\xi(x_3))\in\Lag_n^{3}$ 都是极大的。
\end{itemize}
\end{theorem}

\begin{remark}[11.20]
若 $S_{g,n}$ 是亏格为 $g$、有 $n$ 个边界分支的紧曲面，则 $\mathcal{R}(S_{g,n},\Sp_{2n}(\R))$ 是连通空间，所以 $\mathcal{M}(S_{g,n},\Sp_{2n}(\R))$ 不可能是开且闭的。然而，我们可以考虑 $\mathcal{R}(S_{g,n},\Sp_{2n}(\R))$ 的下述子空间：设 $\mathcal{R}^{p}(S_{g,n},\Sp_{2n}(\R))$ 是所有满足下述条件的表示 $\rho$ 构成的空间——对每个同伦于某个边界分支的 $\gamma\in\pi_1(S_{g,n})$，$\rho(\gamma)$ 固定 $\Lag_n$ 中的某个元素。这样的表示称为外周抛物的（peripherally parabolic）。于是 $\mathcal{M}(S_{g,n},\Sp_{2n}(\R))$ 是 $\mathcal{R}^{p}(S_{g,n},\Sp_{2n}(\R))$ 的开且闭的子集。
\end{remark}

% [原文 p.118]

\section{带标架极大表示空间上的坐标}

在本节中，我们引入极大带标架表示空间上的坐标，其方式类似于映入分裂实李群的正带标架表示的 Fock--Goncharov 坐标。

\subsection{横截与极大的带标架表示}

设 $S$ 是纤毛曲面，配有有限体积的、以测地线为边界的完备双曲结构。设 $\tilde S\subseteq\HH^{2}$ 是 $S$ 的万有覆叠，并设 $\tilde P\subseteq\partial_\infty^{\mathrm{red}}\tilde S$ 是 $S$ 的全部穿孔与纤毛在 $\tilde S$ 边界上的提升。

\begin{definition}[11.21]
（拉格朗日）标架是一个映射 $F\colon\tilde P\to\Lag_n$。设 $\rho\colon\pi_1(S)\to\Sp_{2n}(\R)$ 是同态。$\rho$ 的标架是一个 $\pi_1(S)$-等变的标架 $F\colon\tilde P\to\Lag_n$，即对每个 $\gamma\in\pi_1(S)$ 及一切 $\tilde p\in\tilde P$，有 $F(\gamma(\tilde p))=\rho(\gamma)F(\tilde p)$。带标架同态是一对 $(\rho,F)$，其中 $F$ 是 $\rho$ 的标架。
\end{definition}

所有带标架同态构成的空间记作 $\Hom^{\mathrm{fr}}(\pi_1(S),\Sp_{2n}(\R))$。空间
\[
\mathcal{E}^{\mathrm{fr}}(S,\Sp_{2n}(\R)):=\Hom^{\mathrm{fr}}(\pi_1(S),\Sp_{2n}(\R))/\Sp_{2n}(\R)
\]
称为带标架表示的模空间。带标架表示是 $\mathcal{E}^{\mathrm{fr}}(S,\Sp_{2n}(\R))$ 中的一个元素。

\begin{remark}[11.22]
注意，并非每个同态 $\rho\colon\pi_1(S)\to\Sp_{2n}(\R)$ 都容许标架。容许标架的同态称为外周抛物的。
\end{remark}

设 $\Delta$ 是 $S$ 的一个理想三角剖分，$\tilde\Delta$ 表示 $\Delta$ 到万有覆叠 $\tilde S$ 的提升。

\begin{definition}[11.23]
带标架同态 $(\rho,F)$ 称为 $\Delta$-横截的，如果对任意两个由 $\tilde\Delta$ 的一条边相连的穿孔 $\tilde p_1,\tilde p_2\in\tilde P$，拉格朗日子空间 $F(\tilde p_1)$ 与 $F(\tilde p_2)$ 都是横截的。
\end{definition}

所有 $\Delta$-横截的带标架同态构成的空间记作 $\Hom^{\mathrm{fr}}_{\Delta}(\pi_1(S),\Sp_{2n}(\R))$。空间
\[
\mathcal{E}^{\mathrm{fr}}_{\Delta}(S,\Sp_{2n}(\R)):=\Hom^{\mathrm{fr}}_{\Delta}(\pi_1(S),\Sp_{2n}(\R))/\Sp_{2n}(\R)
\]
称为 $\Delta$-横截的带标架表示的模空间。

\begin{definition}[11.24]
标架 $F$ 称为极大的，如果对每个循环定向的三元组 $(\tilde p_1,\tilde p_2,\tilde p_3)\in\tilde P^{3}$，拉格朗日三元组 $(F(\tilde p_1),F(\tilde p_2),F(\tilde p_3))$ 都是极大的。同态 $\rho\in\Hom(\pi_1(S),\Sp_{2n}(\R))$ 称为极大的，如果它容许一个极大标架 $F$。称对 $(\rho,F)$ 为极大带标架同态。
\end{definition}

% [原文 p.119]

所有极大带标架同态构成的空间记作 $\Hom^{\mathrm{fr}}_{max}(\pi_1(S),\Sp_{2n}(\R))$。空间
\[
\mathcal{M}^{\mathrm{fr}}(S,\Sp_{2n}(\R)):=\Hom^{\mathrm{fr}}_{max}(\pi_1(S),\Sp_{2n}(\R))/\Sp_{2n}(\R)
\]
称为极大带标架表示的模空间。

\begin{remark}[11.25]
由这个定义立即看出，极大带标架表示关于 $S$ 的任意理想三角剖分都是横截的。
\end{remark}

\begin{remark}[11.26]
注意，一般说来，$\Sp_{2n}(\R)$ 在 $\Hom(\pi_1(S),\Sp_{2n}(\R))$ 上以及在 $\Hom^{\mathrm{fr}}(\pi_1(S),\Sp_{2n}(\R))$ 上的作用的轨道并不都是闭的。所以商空间 $\mathcal{E}(S,\Sp_{2n}(\R))$ 与 $\mathcal{E}^{\mathrm{fr}}(S,\Sp_{2n}(\R))$ 都不是 Hausdorff 的。然而，$\mathcal{E}^{\mathrm{fr}}(S,\Sp_{2n}(\R))$ 内部由极大带标架表示构成的空间 $\mathcal{M}^{\mathrm{fr}}(S,\Sp_{2n}(\R))$ 构成一个 Hausdorff 子空间。

此外，忘掉标架的自然映射 $\mathcal{E}^{\mathrm{fr}}(S,\Sp_{2n}(\R))\to\mathcal{E}(S,\Sp_{2n}(\R))$ 将极大带标架表示的空间 $\mathcal{M}^{\mathrm{fr}}(S,\Sp_{2n}(\R))$ 以有限纤维满射到 11.1.5 节讨论过的极大表示空间 $\mathcal{M}(S,\Sp_{2n}(\R))$ 上。这意味着，每个极大表示都容许标架，并且至多有有限多个标架。注意，一般情形下这是不对的。既有不容许标架的表示，也有容许无穷多个标架的表示（更多细节参见 [14]）。
\end{remark}

事实上，要检验标架是否是极大的，只需对某一个三角剖分的所有正定向三角形检验，而不必对 $\tilde P$ 中一切可能的循环定向三元组检验。

\begin{proposition}[11.27]（{[14]} 的 4.12 节）
设 $\Delta$ 是 $S$ 的一个三角剖分，$F$ 是标架，使得对 $\tilde\Delta$ 的每个正定向三角形 $(\tilde p_1,\tilde p_2,\tilde p_3)$，拉格朗日三元组 $(F(\tilde p_1),F(\tilde p_2),F(\tilde p_3))$ 都是极大的。那么标架 $F$ 是极大标架。
\end{proposition}

\subsection{图上的带标架平凡局部系统}

正如我们在 11.1 节开头提到的，管型埃尔米特李群具有非交换代数上辛群的结构。因此读者可以预期，极大表示的坐标将类似于 $\PGL_2(\R)$ 情形下正表示的坐标，只不过取值于某个适当的非交换代数。然而，这只在某种程度上成立，因为非交换代数中有一个额外的复杂之处：代数的可逆元群通过共轭非平凡地作用在代数上。这使得我们无法像对 Fock--Goncharov 坐标那样，对每条内部边、每个理想三角形局部地定义极大表示的坐标。为了顾及这个非平凡的共轭作用，我们需要在一个三角形中作出某种选择，然后把这个选择传播到整个曲面上。此外，还存在一些

% [原文 p.120]

在 $\PGL_2(\R)$ 情形中看不到的额外参数，它们可以看作曲面上一些曲线周围的整体扭转，因而也不是局部的。

为了参数化极大带标架表示的空间，我们将采用如下策略：在曲面 $S$ 上构造一个适应于给定理想三角剖分 $\Delta$ 的图 $\Gamma$；然后对给定的极大带标架表示，在这个图上构造一个带标架局部系统；最后利用这个局部系统给出参数。

\subsubsection{图 $\Gamma$}

我们在 $S$ 上构造如下图 $\Gamma$。$\Gamma$ 的顶点集是 $V_\Gamma:=\{(\tau,r)\in T\times E\mid r\subset\bar\tau\}$。顶点 $(\tau,r)\in V_\Gamma$ 可以看作三角形 $\tau\in T$ 中紧挨着内部边 $r\in E$ 的一个点。

\Fig[0.36\linewidth]{figures/fig_11_1.png}{11.1}{对 $v\in V_\Gamma$，$v^{t}$、$v^{b}$、$v^{r}$ 与 $v^{l}$ 的定义。}

设 $v=(\tau,r)\in V_\Gamma$，并设边 $r$ 连接两个穿孔 $p,p'\in P$。若在 $\tau$ 内用两条互不相交的简单弧段把 $v$ 与 $p$、$p'$ 连接起来，就得到 $\tau$ 内一个以 $v,p,p'$ 为顶点的三角形。若三角形 $(v,p,p')$ 的定向与曲面的定向一致，我们记 $v^{t}:=p$、$v^{b}:=p'$（见图 11.1）。进一步，我们定义 $v^{r}$ 为 $P$ 中唯一使得 $(v^{t},v^{b},v^{r})=\tau$ 的穿孔。若 $r$ 是 $\Delta$ 的内部边，我们还定义 $v^{l}$ 为 $P$ 中唯一使得 $(v^{b},v^{t},v^{l})$ 是 $\Delta$ 的三角形的穿孔。这个三角形沿 $r$ 与 $\tau$ 相邻。

现在我们来描述 $\Gamma$ 的定向边集 $E_\Gamma^{+}$。先固定如下记号：从顶点 $v$ 到顶点 $v'$ 的定向边记作 $(v\to v')$ 与 $(v'\leftarrow v)$。

\begin{itemize}
\item 设 $v=(\tau_1,r)$、$v'=(\tau_2,r)$ 且 $\tau_1<\tau_2$，则 $(v\to v')$ 是 $\Gamma$ 的一条边。它穿过三角剖分 $\Delta$ 的内部边 $r$。
\item 设 $v_1$ 与 $v_2$ 是位于同一三角形中的两个顶点，且 $v_1^{b}=v_2^{t}$，则 $(v_1\to v_2)$ 是 $\Gamma$ 的一条边（见图 11.2）。
\end{itemize}

\begin{remark}[11.28]
图 $\Gamma$ 没有重边，也没有二元环。
\end{remark}

% [Tutor Note: 顶点集的计数自查：V_Γ = {(τ,r) : r ⊂ τ̄} 是“三角形—边”的关联对，故对由 |T| 个理想三角形组成的剖分有 |V_Γ| = 3|T|（每个三角形恰好有三条边）。Fig. 11.2 的四边形剖分成两个三角形，图恰有 6 个顶点 v_1,…,v_6，正合此数；其中 v_1,v_2 紧挨内部边（对应 (τ,r)、(τ′,r) 两个顶点），其余四个紧挨外部边。可用于核对“每条边的两个顶点分属两个相邻三角形”的构造。]

% [原文 p.121]

\Fig[0.62\linewidth]{figures/fig_11_2.png}{11.2}{图 $\Gamma$ 的一个例子。这里三角剖分的右上边是外部边，其余各边都是内部边，且 $\tau<\tau'$。}

我们还记 $E_\Gamma^{-}:=\{(v'\leftarrow v)\mid (v\leftarrow v')\in E_\Gamma^{+}\}$。由于 $\Gamma$ 没有二元环，$E_\Gamma^{+}\cap E_\Gamma^{-}=\varnothing$。我们记 $E_\Gamma:=E_\Gamma^{+}\cup E_\Gamma^{-}$。

我们用 $\tilde\Gamma$ 表示 $\Gamma$ 到 $S$ 的万有覆叠 $\tilde S\subset\HH^{2}$ 的提升。我们还选取一点 $b\in V_\Gamma$ 并固定它的一个提升 $\tilde b\in V_{\tilde\Gamma}$。假定对某 $\tau_0\in T$ 与 $\tilde\tau_0\in\tilde T$，有 $b\in\tau_0$ 且 $\tilde b\in\tilde\tau_0$。

\begin{remark}[11.29]
$S$ 上任何始于并终于 $V_\Gamma$ 的道路都可以形变为 $\Gamma$ 中的一条无向道路。
\end{remark}

为简化记号，到本章末尾我们将记 $G=\Sp_{2n}(\R)$。

\begin{definition}[11.30]
$\tilde\Gamma$ 的横截标架是一对映射 $F^{t},F^{b}\colon V_{\tilde\Gamma}\to\Lag_n$，使得对每个 $v\in V_{\tilde\Gamma}$，$F^{t}(v)$ 与 $F^{b}(v)$ 都横截。
\end{definition}

\begin{definition}[11.31]
$\tilde\Gamma$ 上的平凡 $G$-局部系统是一个映射 $T\colon V_{\tilde\Gamma}\cup E_{\tilde\Gamma}\to G$，使得对任意边 $(v'\leftarrow v)\in E_{\tilde\Gamma}$（$v,v'\in V_{\tilde\Gamma}$），有 $T(v'\leftarrow v)=T(v\leftarrow v')^{-1}$ 以及 $T(v')=T(v'\leftarrow v)T(v)$。由映射 $T$ 定义的 $\tilde\Gamma$ 上的平凡 $G$-局部系统记作 $(\tilde\Gamma,T)$。
\end{definition}

\begin{remark}[11.32]
由定义 11.31 可知，对每个环 $(e_1,\dots,e_k)$（其中一切 $e_i\in E_{\tilde\Gamma}$）都有 $T(e_k)\cdots T(e_1)=1$，即绕环的和乐是平凡的。这正是把这样的局部系统称为平凡的原因。
\end{remark}

设 $v\in V_{\tilde\Gamma}$ 是 $\tilde\Gamma$ 的一个固定顶点。映射 $T$ 由 $T|_{E_{\tilde\Gamma}^{+}}$ 与 $T(v)$ 决定。元素 $T(v)$ 以及对一切 $e\in E_{\tilde\Gamma}^{+}$ 的元素 $T(e)$ 可以是 $G$ 的任意元素，只要满足下述条件：对每个环 $(e_1,\dots,e_k)$（其中一切 $e_i\in E_{\tilde\Gamma}$）都有 $T(e_k)\cdots T(e_1)=1$。

设 $(F^{t},F^{b})$ 是横截标架，$(\tilde\Gamma,T)$ 是 $\tilde\Gamma$ 上的平凡 $G$-局部系统。我们说横截标架 $(F^{t},F^{b})$ 适应（adapted）于平凡局部系统 $(\tilde\Gamma,T)$，如果对每个 $v\in V_{\tilde\Gamma}$ 都有 $T(v)F^{t}(v)=L_0^{\mathrm{opp}}$ 且 $T(v)F^{b}(v)=L_0$。四元组 $\mathfrak{L}:=(\tilde\Gamma,T,F^{t},F^{b})$ 称为横截带标架平凡 $G$-局部系统，如果 $(F^{t},F^{b})$ 是适应于 $(\tilde\Gamma,T)$ 的横截标架。

% [原文 p.122]

由于 $\tilde\Gamma$ 没有重边，$\tilde\Gamma$ 中的任何道路 $\gamma$ 都可以唯一地写成顶点序列 $\gamma=(v_k\leftarrow v_{k-1}\leftarrow\cdots\leftarrow v_1)$，其中 $v_1$ 是 $\gamma$ 的起点、$v_k$ 是 $\gamma$ 的终点，且对一切 $i\in\{1,\dots,k-1\}$ 有 $(v_{i+1}\leftarrow v_i)\in E_{\tilde\Gamma}$。我们如下把映射 $T$ 延拓到道路空间上：若 $\gamma=(v_k\leftarrow v_{k-1}\leftarrow\cdots\leftarrow v_1)$，则 $T(\gamma):=T(v_k\leftarrow v_{k-1})\cdots T(v_2\leftarrow v_1)$。

\begin{remark}[11.33]
设 $\gamma$ 是 $\tilde\Gamma$ 中连接两个顶点 $v$ 与 $w$ 的道路。那么实际上 $T(\gamma)$ 只依赖于 $v$ 与 $w$，而不依赖于整条道路 $\gamma$。因此，为了强调这一点，有时我们将直接写 $T(w\leftarrow v)$ 而不写 $T(\gamma)$。
\end{remark}

\subsubsection{从局部系统到表示}

设 $\mathfrak{L}=(\tilde\Gamma,T,F^{t},F^{b})$ 是 $\tilde\Gamma$ 上的横截带标架平凡 $G$-局部系统。基本群 $\pi_1(S)$ 在 $\tilde S$ 上的作用限制为在 $\tilde\Gamma$ 上的作用。

\begin{definition}[11.34]
局部系统 $\mathfrak{L}=(\tilde\Gamma,T,F^{t},F^{b})$ 称为 $\pi_1(S)$-不变的，如果对一切 $\gamma\in\pi_1(S)$ 与一切边 $(w\leftarrow v)\in E_{\tilde\Gamma}$（其中 $v,w\in V_{\tilde\Gamma}$），有 $T(w\leftarrow v)=T(\gamma w\leftarrow\gamma v)$。
\end{definition}

$\pi_1(S)$-不变的 $G$-局部系统总以如下方式给出一个同态 $\rho\colon\pi_1(S)\to G$：取 $v\in V_{\tilde\Gamma}$，对每个 $\gamma\in\pi_1(S)$ 定义
\begin{equation}
\rho_{\mathfrak{L}}(\gamma):=T(\gamma v)^{-1}T(v). \tag{11.2}
\end{equation}

\begin{remark}[11.35]
若 $\mathfrak{L}=(\tilde\Gamma,T,F^{t},F^{b})$ 是 $\pi_1(S)$-不变的局部系统，则标架 $(F^{t},F^{b})$ 是 $\rho_{\mathfrak{L}}$-等变的，即对一切 $v\in V_{\tilde\Gamma}$，$F_t(\gamma v)=\rho_{\mathfrak{L}}(\gamma)(F_t(v))$ 且 $F_b(\gamma v)=\rho_{\mathfrak{L}}(\gamma)(F_b(v))$。
\end{remark}

\begin{proposition}[11.36]
映射 $\rho_{\mathfrak{L}}$ 是同态，且不依赖于 $v\in V_{\tilde\Gamma}$ 的选取。
\end{proposition}

\begin{proof}
设 $w\in V_{\tilde\Gamma}$ 是另一个顶点。则
\begin{align*}
T(\gamma w)^{-1}T(w)&=\left(T(\gamma w\leftarrow\gamma v)T(\gamma v)\right)^{-1}\left(T(w\leftarrow v)T(v)\right)\\
&=T(\gamma v)^{-1}T(\gamma w\leftarrow\gamma v)^{-1}T(w\leftarrow v)T(v)\\
&=T(\gamma v)^{-1}T(v),
\end{align*}
因为由 $H$-不变性有 $T(\gamma w\leftarrow\gamma v)=T(w\leftarrow v)$。所以 $\rho_{\mathfrak{L}}$ 不依赖于 $v\in V_{\tilde\Gamma}$ 的选取。

现在设 $\gamma_1,\gamma_2\in\pi_1(S)$。我们得到 $\rho_{\mathfrak{L}}(\gamma_2)=T(\gamma_2(\gamma_1 v))^{-1}T(\gamma_1 v)$ 以及 $\rho_{\mathfrak{L}}(\gamma_1)=T(\gamma_1 v)^{-1}T(v)$。因此
\begin{align*}
\rho_{\mathfrak{L}}(\gamma_2)\rho_{\mathfrak{L}}(\gamma_1)
&=T(\gamma_2(\gamma_1 v))^{-1}T(\gamma_1 v)T(\gamma_1 v)^{-1}T(v)\\
&=T(\gamma_2(\gamma_1 v))^{-1}T(v)=\rho_{\mathfrak{L}}(\gamma_2\gamma_1).
\end{align*}
所以 $\rho_{\mathfrak{L}}$ 是同态。
\end{proof}

% [Tutor Note: 记号提醒（易错点）：本节大部分地方用上标 F^t, F^b（Def. 11.30/11.31），而 Remark 11.35 与 11.36 的证明却用下标 F_t, F_b，指的是同一对映射，原文如此，译文照印。另：11.36 的证明中出现“由 H-不变性”，但本章从未定义 H；这里实际使用的是 Def. 11.34 的 π_1(S)-不变性（即 T(γw ← γv) = T(w ← v)，正是 Def. 11.34 的等式），读作 π_1(S)-不变性即可。]

% [原文 p.123]

每个标架 $F\colon\tilde P\to\Lag_n$ 都如下定义 $\tilde\Gamma$ 上唯一的标架 $(F^{t},F^{b})$：对 $\tilde\Gamma$ 的每个顶点 $v$，定义 $F^{t}(v):=F(v^{t})$ 与 $F^{b}(v):=F(v^{b})$（见图 11.1）。此时我们说 $(F^{t},F^{b})$ 由 $F$ 诱导（induced）。我们还记 $F^{r}(v):=F(v^{r})$ 与 $F^{l}(v):=F(v^{l})$。

设 $r$ 是 $\tilde\Gamma$ 的一条内部边，它分隔两个三角形 $\tau$ 与 $\tau'$，并设 $v:=(\tau,r)$、$v':=(\tau',r)$。则 $F^{t}(v')=F^{b}(v)$、$F^{b}(v')=F^{t}(v)$、$F^{r}(v')=F^{l}(v)$，且 $F^{l}(v')=F^{r}(v)$。进一步，设 $v$ 与 $v'$ 是 $\tilde\Gamma$ 中位于同一三角形内、且 $\tilde\Gamma$ 中有边 $(v\to v')$ 的两个顶点。则 $F^{t}(v')=F^{b}(v)$、$F^{b}(v')=F^{r}(v)$、$F^{r}(v')=F^{t}(v)$。

下面的命题是直接的：

\begin{proposition}[11.37]
设 $(F^{t},F^{b})$ 是 $\tilde\Gamma$ 上的横截标架，使得
\begin{itemize}
\item 对 $\tilde\Gamma$ 的每条分隔两个三角形 $\tau$ 与 $\tau'$ 的内部边 $r$，若 $v:=(\tau,r)$、$v':=(\tau',r)$，则 $F^{t}(v')=F^{b}(v)$；
\item 对 $\tilde\Gamma$ 中位于同一三角形内、且 $\tilde\Gamma$ 中有边 $(v\to v')$ 的任意两个顶点 $v,v'$，有 $F^{t}(v')=F^{b}(v)$。
\end{itemize}
那么存在唯一的标架 $F\colon\tilde P\to\Lag_n$ 诱导出 $(F^{t},F^{b})$。
\end{proposition}

\begin{remark}[11.38]
设 $(\rho,F)$ 是带标架同态。若 $(F^{t},F^{b})$ 是由 $F$ 诱导的标架，则 $(F^{t},F^{b})$ 也是 $\rho$-等变的。若 $\mathfrak{L}=(\tilde\Gamma,T,F^{t},F^{b})$ 是 $\tilde\Gamma$ 上 $\pi_1(S)$-不变的横截带标架平凡 $G$-局部系统，则 $F$ 是 $\rho_{\mathfrak{L}}$ 的标架，但一般说来 $\rho\ne\rho_{\mathfrak{L}}$。
\end{remark}

\begin{definition}[11.39]
我们说横截带标架平凡 $G$-局部系统 $(\tilde\Gamma,T,F^{t},F^{b})$ 是极大的，如果标架 $(F^{t},F^{b})$ 由某个极大标架 $F\colon\tilde P\to\Lag_n$ 诱导。
\end{definition}

\subsubsection{从表示到局部系统}

本节中，我们说明如何由极大带标架表示出发，在 $\Gamma$ 上构造一个极大带标架平凡 $G$-局部系统。设 $\Delta$ 是 $S$ 的一个理想三角剖分，$\tilde\Gamma$ 是按 11.2.2.1 节定义的 $\tilde S$ 上的图。给定极大带标架同态 $(\rho,F)\in\Hom^{\mathrm{fr}}_{+}(\pi_1(S,b),G)$，标架 $F$ 给出 $\tilde\Gamma$ 的 $\rho$-等变横截标架 $(F^{t},F^{b})$。我们将定义一个局部系统 $\mathfrak{L}=(\tilde\Gamma,T)$，使得 $(F^{t},F^{b})$ 成为它的横截标架，且 $\rho_{\mathfrak{L}}=\rho$。

我们先从 $S$ 是三角形以及剖分四边形这两个最简单的例子开始，然后描述对一般曲面 $S$ 的步骤。

\textbf{三角形}

设 $S=\tilde S$ 是顶点为 $p_1,p_2,p_3$ 的理想三角形。$S$ 的唯一的三角剖分 $\Delta$ 由 $S$ 的全部外部边构成。设 $(\rho,F)$ 是极大带标架同态。由于 $\pi_1(S)$ 平凡，$\rho$ 也平凡。设 $F(p_i)=F_i$，$i\in\{1,2,3\}$。那么 $(F_1,F_2,F_3)$ 是极大三元组。设图 $\Gamma=\tilde\Gamma$ 由三个顶点 $v$、$v'$ 与 $v''$ 以及三条边 $(v'\leftarrow v)$、$(v''\leftarrow v')$ 与 $(v\leftarrow v'')$ 构成。

% [原文 p.124]

由极大性，存在元素 $g\in G$，使得 $g(F_1,F_2,F_3)=(L_0^{\mathrm{opp}},L_0,L^{+})$（见图 11.3）。

\Fig[0.42\linewidth]{figures/fig_11_3.png}{11.3}{三角形中的图 $\Gamma$。}

我们定义 $T(v):=g$ 以及 $T(v'\leftarrow v):=T(v''\leftarrow v'):=T(v\leftarrow v''):=\begin{pmatrix} -1 & 1\\ -1 & 0\end{pmatrix}$，于是 $T(v')=T(v'\leftarrow v)T(v)$、$T(v'')=T(v''\leftarrow v')T(v')$，并且
\begin{align*}
T(v')F^{t}(v')&=T(v')F_2=T(v'\leftarrow v)L_0=L_0^{\mathrm{opp}},\\
T(v')F^{b}(v')&=T(v')F_3=T(v'\leftarrow v)L^{+}=L_0,\\
T(v')F^{r}(v')&=T(v')F_1=T(v'\leftarrow v)L_0^{\mathrm{opp}}=L^{+}.
\end{align*}
类似地可以看出 $T(v'')F^{t}(v'')=L_0^{\mathrm{opp}}$、$T(v'')F^{b}(v'')=L_0$ 以及 $T(v'')F^{r}(v'')=L^{+}$。此外，$T(v\leftarrow v'')T(v''\leftarrow v')T(v'\leftarrow v)=\mathrm{Id}$。于是我们在 $\Gamma$ 中构造出了一个定义良好的极大带标架平凡 $G$-局部系统。注意，这个局部系统并不唯一，它依赖于 $g\in G$ 的选取。这个 $g$ 在左乘 $\mathrm{O}(n)$ 中元素的意义下唯一确定（见命题 11.14）。

\textbf{四边形。跨过三角剖分的一条边}

设 $S=\tilde S$ 是顶点按循环次序为 $p_1,p_3,p_2,p_4$ 的理想四边形。设 $\Delta$ 是由 $S$ 的全部外部边以及一条连接 $p_1$ 与 $p_2$ 的内部边构成的三角剖分。设 $(\rho,F)$ 是极大带标架同态。由于 $\pi_1(S)$ 平凡，$\rho$ 也平凡。设 $F(p_i)=F_i$，$i\in\{1,\dots,4\}$。那么 $(F_1,F_4,F_2,F_3)$ 是正四元组。图 $\Gamma=\tilde\Gamma$ 有六个顶点，其中四个紧挨外部边，另两个 $v$ 与 $v'$ 紧挨内部边。于是存在边 $(v'\leftarrow v)$ 或 $(v\leftarrow v')$。不失一般性，我们假定是前一种情形。

% [原文 p.125]

\Fig[0.92\linewidth]{figures/fig_11_4.png}{11.4}{跨过三角剖分的一条边。}

由命题 11.16 给出的横截性，存在元素 $g\in G$，使得
\[
g(F_2,F_3,F_1,F_4)=(L_0,L^{+},L_0^{\mathrm{opp}},L(a)),
\]
其中 $L(a)$ 由基 $\mathbf{f}_0-\mathbf{e}_0 a$ 生成，$a\in\Sym_n^{+}(\R)$（见图 11.4）。对每个矩阵 $a\in\Sym_n^{+}(\R)$，存在唯一的平方根 $q\in\Sym_n^{+}(\R)$，使得 $q^{2}=a$。我们定义 $T(v):=g$、$T(v'\leftarrow v):=\begin{pmatrix} 0 & \sqrt{a^{-1}}\\ -\sqrt{a} & 0\end{pmatrix}$，于是 $T(v')=T(v'\leftarrow v)T(v)$。$\Gamma$ 的其余各边都落在 $S$ 的两个三角形之一内部，所以我们如上一小节那样为它们构造映射 $T$。这样我们便在 $\Gamma$ 上构造出了一个极大带标架平凡 $G$-局部系统。元素 $g$ 在左乘 $\mathrm{O}(n)$ 中元素的意义下唯一。

现在设 $S$ 是任意纤毛曲面，$\Delta$ 是 $S$ 的理想三角剖分，$(\rho,F)\in\Hom^{\mathrm{fr}}_{+}(\pi_1(S),G)$。设 $\tilde\Gamma$ 是 11.2.2.1 节中的图，其标架 $(F^{t},F^{b})$ 由 $F$ 诱导。我们可以假定基本群的基点与某个位于三角形 $\tau$ 中的点 $b\in V_\Gamma$ 一致。取 $b$ 的提升 $\tilde b\in V_{\tilde\Gamma}$。

从 $\tilde b$ 出发，对每个三角形反复运用 11.2.2.3 节描述的步骤、对每对相邻三角形反复运用 11.2.2.3 节的步骤，我们便得到 $\tilde\Gamma$ 上的一个极大带标架局部系统。这个局部系统具有下述性质：(11.2) 中描述的表示与 $\rho$ 一致。然而，这个局部系统一般不是 $\pi_1(S)$-不变的，尽管标架是 $\rho$-等变的。在 11.2.3 节定理 11.40 的证明中，我们将修改这个局部系统，使之成为 $\pi_1(S)$-不变的。

\subsection{极大表示空间的参数化}

现在我们准备给出极大带标架表示空间的一个参数化。和前面一样，设 $S$ 是欧拉示性数为 $\chi(S)$、有 $p_e$ 个纤毛的纤毛曲面。

% [原文 p.126]

\begin{theorem}[11.40]
空间 $\mathcal{M}^{\mathrm{fr}}(S,G)$ 同胚于下述空间：
\[
\left(\Sym_n^{+}(\R)^{p_e-3\chi(S)}\times\mathrm{O}(n)^{1-\chi(S)}\right)/\mathrm{O}(n),
\]
其中 $\mathrm{O}(n)$ 在每个因子上通过共轭作用。
\end{theorem}

\begin{remark}[11.41]
正如我们在 11.2.2 节开头所宣布的，有些参数确实是非交换代数 $\Mat_n(\R)$ 的元素，具体说是它的“正的部分” $\Sym_n^{+}(\R)$ 的元素。和在 $\PGL_2(\R)$ 情形中一样，它们与三角剖分的内部边相联系。但我们还看到了额外的 $\mathrm{O}(n)$-参数以及按 $\mathrm{O}(n)$-作用取共轭的商，这些在 Fock--Goncharov 情形中是没有的。为说明它们的含义，我们给出这个定理的证明，尽管它相当技术性。
\end{remark}

\begin{proof}
设 $(\rho,F)$ 是极大带标架同态。我们取 $S$ 的某个理想三角剖分 $\Delta$，其内部边集为 $E_{in}$。标架 $F$ 诱导出 $\tilde\Gamma$ 的标架 $(F^{t},F^{b})$。

设 $b\in S$ 是 $\pi_1(S,b)$ 的基点，$\tilde b\in\tilde S$ 是 $b$ 的提升。如前，假定 $\tilde b\in V_{\tilde\Gamma}$。设 $g\in G$ 使得 $g(F^{t}(\tilde b),F^{b}(\tilde b),F^{r}(\tilde b))=(L_0^{\mathrm{opp}},L_0,L^{+})$。现在我们定义
\[
\begin{aligned}
X:=\bigl\{(\rho,F,g)\ \bigm|\ &(\rho,F)\in\Hom^{\mathrm{fr}}_{+}(\pi_1(S,b),G),\ g\in G,\ \text{使得：}\\
&g(F^{t}(\tilde b),F^{b}(\tilde b),F^{r}(\tilde b))=(L_0^{\mathrm{opp}},L_0,L^{+})\bigr\},
\end{aligned}
\]
以及 $X$ 上的等价关系 $\sim$ 如下：
\[
(\rho,F,g)\sim(\rho',F',g'),\quad\text{当且仅当 }(g\rho g^{-1},gF)=(g'\rho'g'^{-1},g'F').
\]
\emph{步骤 1：}假定 $S=\tilde S$ 是三角剖分多边形，$\Gamma$ 按 11.2.2.1 节定义。设 $(\rho,F,g)\in X$。我们定义 $T(b)=g$，然后应用 11.2.2.3 节定义的步骤。这样我们得到 $\Gamma$ 上的一个带标架横截 $G$-局部系统。设 $v=(\tau,r)\in V_\Gamma$，其中 $r\in E_{in}$，则 $v$ 位于以 $v^{t},v^{b},v^{r}\in P$ 为顶点的三角形 $\tau$ 中。设 $\tau'$ 是沿边 $r=(v^{t},v^{b})$ 与 $\tau$ 相邻的三角形，而 $v^{b},v^{t},v^{l}$ 是 $\tau'$ 的顶点。设 $b_v\in\Sym_n^{+}(\R)$ 由 $T(v)F(v^{l})=L(b_v)$ 定义，其中 $L(b_v)$ 是由基 $\mathbf{f}_0-\mathbf{e}_0 b_v$ 生成的拉格朗日子空间。元素 $b_v$ 对一切 $v\in V_\Gamma$ 都唯一确定。而且，由 11.2.2.3 节的构造，$b_v=b_{v'}$，其中 $v'=(\tau',r)$。因此，我们可以把元素 $b_v$ 不联系于 $\Gamma$ 的顶点，而联系于 $\Delta$ 的内部边 $r$。所以我们改记作 $b_r$。

于是我们得到下述映射：
\[
\mathcal{X}\colon \quad X\longrightarrow\Sym_n^{+}(\R)^{E_{in}},\qquad
(\rho,F,g)\longmapsto(b_r)_{r\in E_{in}}.
\]
注意，$\mathcal{X}(\rho,F,g)=\mathcal{X}(\rho',F',g')$ 当且仅当 $(\rho,F,g)\sim(\rho',F',g')$。这个映射是满的，因为任意选取 $(b_e)_{e\in E_{in}}$ 都定义 $\Gamma$ 上的一个横截带标架 $G$-局部系统，它由 11.2.2.2 节给出一个带标架同态 $(\rho,F)\in\Hom^{\mathrm{fr}}_{\Delta}(\pi_1(S),G)$。

% [原文 p.127]

我们还有 $\mathrm{O}(n)$ 在 $X$ 上的作用：它在最后一个分量上通过左乘作用，在 $\Sym_n^{+}(\R)^{E_{in}}$ 上通过逐分量共轭作用，且这个作用与 $\mathcal{X}$ 相容，即：设 $\diag(u,u)=k\in\mathcal{K}$，$u\in\mathrm{O}(n)$，则 $\mathcal{X}(\rho,D,kg)=(ub_r u^{-1})_{r\in E_{in}}$。我们得到商映射
\[
\bar{\mathcal{X}}\colon\mathcal{M}^{\mathrm{fr}}(S,G)\longrightarrow\Sym_n^{+}(\R)^{E_{in}}/\mathrm{O}(n),
\]
它是一个同胚。

\emph{步骤 2：}现在转向一般情形。我们在 $\tilde S$ 中选取 $S$ 的一个连通基本域 $S_0$，它完全由理想三角形组成，且 $\tilde b$ 落在其中之一。基本域 $S_0$ 是一个理想多边形，所以我们可以应用上面的步骤，得到 $\Gamma_0$ 上的一个带标架横截 $G$-局部系统，其中 $\Gamma_0$ 是 $\tilde\Gamma$ 的含于 $S_0$ 的子图。我们希望以 $\pi_1(S)$-不变的方式把这个局部系统延拓到整个图 $\tilde\Gamma$ 上。

为得到 $\pi_1(S)$-不变的局部系统，按 11.2.2.2 节，我们需要对每个 $\tilde v\in V_{\tilde\Gamma}$ 如下定义 $T(\tilde v)$：设 $\gamma\in\pi_1(S)$ 是使得 $\gamma v\in S_0$ 的唯一元素。我们定义 $T(\tilde v):=T(\gamma\tilde v)\rho(\gamma)$。由此立即得到：对每个 $\tilde v\in V_{\tilde\Gamma}$ 及每个 $\gamma\in\pi_1(S)$，$T(\gamma\tilde v)=T(\tilde v)\rho(\gamma)^{-1}$。

局部系统 $(\tilde\Gamma,T)$ 确实是 $\pi_1(S)$-不变的，因为若 $(w\leftarrow v)\in E_{\tilde\Gamma}$ 且 $\gamma\in\pi_1(S)$，则
\begin{align*}
T(\gamma\tilde w\leftarrow\gamma\tilde v)&=T(\gamma\tilde w)T(\gamma\tilde v)^{-1}\\
&=T(\tilde w)\rho(\gamma)^{-1}\rho(\gamma)T(\tilde v)^{-1}\\
&=T(\tilde w)T(\tilde v)^{-1}\\
&=T(\tilde w\leftarrow\tilde v).
\end{align*}
此外，$(\tilde\Gamma,T)$ 的标架 $(F^{t},F^{b})$ 适应于 $(\tilde\Gamma,T)$。确实，若 $\tilde v$ 含于 $S_0$，则由构造有 $T(\tilde v)(F^{t}(\tilde v),F^{b}(\tilde v))=(L_0^{\mathrm{opp}},L_0)$。若 $\tilde v\in V_{\tilde\Gamma}$ 是任一顶点，则存在唯一的 $\gamma\in\pi_1(S)$ 使得 $\gamma\tilde v$ 含于 $S_0$。于是
\begin{align*}
T(\tilde v)(F^{t}(\tilde v),F^{b}(\tilde v))
&=T(\gamma\tilde v)\rho(\gamma)\rho(\gamma)^{-1}(F^{t}(\gamma\tilde v),F^{b}(\gamma\tilde v))\\
&=T(\gamma\tilde v)(F^{t}(\gamma\tilde v),F^{b}(\gamma\tilde v))\\
&=(L_0^{\mathrm{opp}},L_0).
\end{align*}
带标架局部系统 $(\tilde\Gamma,T,F^{t},F^{b})$ 由它在 $\Gamma_0$ 上的限制以及映射 $T(\tilde w\leftarrow\tilde v)$（$\tilde v\in S_0$、$\tilde w\notin S_0$）决定。因此我们需要理解映射 $T(\tilde w\leftarrow\tilde v)$。

设 $e,e'$ 是三角剖分 $\tilde\Delta$ 的两条边，它们位于 $S_0$ 的边界上但不属于 $\tilde S$ 的边界，并设 $\gamma\in\pi_1(S)$ 是使得 $\gamma(e)=e'$ 的唯一元素。不失一般性，假定 $(\tilde w\leftarrow\tilde v)\in E_{\tilde\Gamma}^{+}$ 穿过 $e$，且 $\tilde v\in S_0$、$\tilde w\notin S_0$。那么边 $(\gamma_i\tilde w\leftarrow\gamma_i\tilde v)\in E_{\tilde\Gamma}^{+}$ 穿过 $e'$，且 $\gamma_i\tilde v\notin S_0$、$\gamma_i\tilde w\in S_0$。此外，
\[
T(\tilde w\leftarrow\tilde v)=T(\tilde w)T(\tilde v)^{-1}=T(\gamma\tilde w)\rho(\gamma)T(\tilde v)^{-1},
\]

% [原文 p.128]

于是
\begin{align*}
T(\tilde w\leftarrow\tilde v)(L_0^{\mathrm{opp}},L_0)
&=T(\gamma\tilde w)\rho(\gamma)T(\tilde v)^{-1}(L_0^{\mathrm{opp}},L_0)\\
&=T(\gamma\tilde w)\rho(\gamma)(F^{t}(\tilde v),F^{b}(\tilde v))\\
&=T(\gamma\tilde w)(F^{t}(\gamma\tilde v),F^{b}(\gamma\tilde v))\\
&=T(\gamma\tilde w)(F^{b}(\gamma\tilde w),F^{t}(\gamma\tilde w))\\
&=(L_0,L_0^{\mathrm{opp}}).
\end{align*}
这意味着 $T(\tilde w\leftarrow\tilde v)=l_e\omega$，其中 $l_e\in L$ 是唯一确定的元素，而 $\omega=\begin{pmatrix} 0 & \mathrm{Id}\\ -\mathrm{Id} & 0\end{pmatrix}$。于是对每一对由 $\pi_1(S)$ 的非平凡元素相联系的外部边 $\{e,e'\}$，我们都得到一个 $L$ 的元素。由极分解，元素 $l_e$ 可以唯一地写成 $l_e=\diag(u_e b_e,u_e b_e^{-1})$，其中 $u_e\in\mathrm{O}(n)$、$b_e\in\Sym_n^{+}(\R)$。我们把如上所有 $\{e,e'\}$ 对构成的集合记作 $E_0$。它的基数是 $1-\chi(S)$。

与步骤 1 一样，我们得到映射
\[
\mathcal{X}\colon \quad X\longrightarrow\Sym_n^{+}(\R)^{E_{in}}\times\mathrm{O}(n)^{E_0},\qquad
(\rho,F,g)\longmapsto(b_r,u_e)_{r\in E_{in},\,e\in E_0}.
\]
注意，$\mathcal{X}(\rho,F,g)=\mathcal{X}(\rho',F',g')$ 当且仅当 $(\rho,F,g)\sim(\rho',F',g')$。进一步，这个映射是满的，因为任意选取 $(b_r,u_e)_{r\in E_{in},\,e\in E_0}$ 都定义 $\tilde\Gamma$ 上一个 $\pi_1(S)$-不变的横截带标架 $G$-局部系统，而它又由 11.2.2.2 节给出一个带标架同态 $(\rho,F)\in\Hom^{\mathrm{fr}}_{max}(\pi_1(S),G)$。

我们还有 $\mathrm{O}(n)$ 在 $X$ 上的作用：它在最后一个分量上通过左乘作用，在 $\Sym_n^{+}(\R)^{E_{in}}\times\mathrm{O}(n)^{E_0}$ 上通过逐分量共轭作用，且这个作用与 $\mathcal{X}$ 相容，即
\[
\mathcal{X}(\rho,D,kg)=(ub_r u^{-1},\,u u_e u^{-1})_{r\in E_{in},\,e\in E_0},
\]
其中 $k=\diag(u,u)$。我们由此得到商映射
\[
\bar{\mathcal{X}}\colon\mathcal{M}^{\mathrm{fr}}(S,G)\longrightarrow\left(\Sym_n^{+}(\R)^{E_{in}}\times\mathrm{O}(n)^{E_0}\right)/\mathrm{O}(n),
\]
它是一个同胚。
\end{proof}

% [Tutor Note: 定理 11.40 的指数自查：对 χ(S) = 1 的理想多边形（p_e 个顶点/穿孔，被剖分成 p_e − 2 个理想三角形）而言，其三角剖分的内部边数恒为 p_e − 3，而 1 − χ(S) = 0 故 E_0 = ∅，于是参数空间为 Sym_n^+(R)^{p_e−3}（再模 O(n) 在此集上的平凡作用）。三角形 p_e = 3 给出 0 个参数（局部系统只由 g 决定，而 g 再模 O(n)），四边形 p_e = 4 给出 1 个参数 a ∈ Sym_n^+(R)，正与本节两个例子的构造相吻合；可用来核对 p_e − 3χ(S) 与 1 − χ(S) 这两个指数。]

% [Tutor Note: 两处印刷瑕疵，照印并说明：① 本页（p.127）与 p.128 中 O(n)-作用的表达式写作 $\mathcal{X}$(ρ, D, kg)，其中 D 未定义，按上下文应为 F（三元组是 (ρ, F, g)）；② p.127 中“边 (γ_i w̃ ← γ_i ṽ)”的下标 i 未定义，按上下文应读作 γ（或 γ 的某个幂）；这些都不影响结论，原文如此。另：p.123、p.125、p.126 处同态空间记号混用 Hom^{fr}_max 与 Hom^{fr}_+（p.165 的段落与 p.126 的 X 定义用后者），二者指同一空间。]

\begin{corollary}[11.42]
空间 $\mathcal{M}^{\mathrm{fr}}(S,G)$ 同伦等价于 $\mathrm{O}(n)^{1-\chi(S)}/\mathrm{O}(n)$，并且有 $2^{1-\chi(S)}$ 个连通分支。
\end{corollary}

关于这个结果的证明，我们请读者参阅 [14, 第 7 节]。
